\subsection{\texorpdfstring{通解$=$特解$+$齐次解}{通解=特解+齐次解}}
前一小节中对解集的种种描述都符合同一个模式：
它们都是一个向量——该方程组的一个特解——
加上另外一些向量的无限制组合。
\nearbyexample{ex:ManyParamsInfManySolsSystem} 给出的解集
例示了这一点。
\begin{equation*}
  \set{
   \underbracket[.7pt]{
     \colvec[r]{0 \\ 4 \\ 0 \\ 0 \\ 0}}_{\text{\shortstack{\rule{0pt}{2ex}特解}}}+
   \underbracket[.7pt]{w\colvec[r]{1 \\ -1 \\ 3 \\ 1 \\ 0}+
       u\colvec[r]{1/2 \\ -1 \\ 1/2 \\ 0 \\ 1}}_{\text{\shortstack{\rule{0pt}{2ex}无限制\\
                                                                组合}}}
       \suchthat w,u\in\Re}
\end{equation*}
组合之所以是无限制的，是指 $w$ 和 $u$ 可以取任何实数\Dash 没有
类似“要求满足 $2w-u=0\mspace{1mu}$”这样的条件来限制
我们可以使用哪些数对 $w,u$。

前面的例题展示了一个符合该模式的无穷解集。
另外两类解集也同样符合该模式。
单元素解集符合该模式，是因为它有一个特解，
而无限制组合那一部分是平凡的。
也就是说，它不是两个向量的组合，也不是一个向量的组合，
而是零个向量的组合。
（按照约定，空向量集的和是零向量。）
空解集符合该模式，是因为此时没有特解，
从而不存在这种形式的和。

\begin{theorem} \label{th:GenEqPartPlusHomo}
%<*th:GenEqPartPlusHomo>
任一线性方程组的
解集都有形式
\begin{equation*}
   \set{\vec{p}+c_1\vec{\beta}_1+\,\cdots\,+c_k\vec{\beta}_k
     \suchthat c_1,\,\ldots\,,c_k\in\Re}
\end{equation*}
其中 \( \vec{p} \) 是任意一个特解，
而向量 $\vec{\beta}_1$, \ldots, $\vec{\beta}_k$ 的个数
等于该方程组经高斯消元后自由变量的个数。
%</th:GenEqPartPlusHomo>
\end{theorem}

解的描述分为两部分：
特解 $\vec{p}$，
以及诸 $\vec{\beta}$ 的无限制线性组合。
我们将用与之相应的两个引理来证明该定理。

我们首先关注无限制的组合。
为此，我们考虑以零向量
作为特解的方程组，
这样便能将 $\vec{p}+c_1\vec{\beta}_1+\dots+c_k\vec{\beta}_k$
缩短为 $c_1\vec{\beta}_1+\dots+c_k\vec{\beta}_k$。

\begin{definition} \label{df:HomogeneousEquation}
%<*df:HomogeneousEquation>
如果一个线性方程带有常数零，从而可以写成
$a_1x_1+a_2x_2+\,\cdots\,+a_nx_n=0$ 的形式，
则称它是\definend{齐次的}\index{homogeneous equation}%
\index{linear equation!homogeneous}。
%</df:HomogeneousEquation>
\end{definition}

\begin{example}  \label{ex:FirstExHomoSys}
对任一线性方程组，比如
\begin{equation*}
  \begin{linsys}{2}
    3x  &+  &4y  &=  3  \\
    2x  &-  &y   &=  1
  \end{linsys}
\end{equation*}
我们将右端设为零，从而得到与之相应的
一个齐次方程组。
\begin{equation*}
  \begin{linsys}{2}
    3x  &+  &4y  &=  0  \\
    2x  &-  &y   &=  0
  \end{linsys}
\end{equation*}
把原方程组的化简
\begin{equation*}
  \begin{linsys}{2}
    3x  &+  &4y  &=  3  \\
    2x  &-  &y   &=  1
  \end{linsys}
  \grstep{-(2/3)\rho_1+\rho_2}
  \begin{linsys}{2}
    3x  &+  &4y        &=  3  \\
        &   &-(11/3)y   &=  -1
   \end{linsys}
\end{equation*}
与相应的齐次方程组的化简加以比较。
\begin{equation*}
  \begin{linsys}{2}
    3x  &+  &4y  &=  0  \\
    2x  &-  &y   &=  0
  \end{linsys}
  \grstep{-(2/3)\rho_1+\rho_2}
  \begin{linsys}{2}
    3x  &+  &4y        &=  0  \\
        &   &-(11/3)y   &=  0
   \end{linsys}
\end{equation*}
显然这两个化简的进行方式完全相同。
于是我们可以通过研究如何化简相应的齐次方程组，
来代替研究如何化简一个线性方程组。
\end{example}

研究相应的齐次方程组比研究原方程组
有一个很大的优点。
非齐次方程组可能无解。
但齐次方程组必定相容，因为它总是至少有
一个解：零向量。


\begin{example} \label{ex:HomoZeroOnlySol}
有些齐次方程组以零向量作为它们唯一的解。
\begin{equation*}
   \begin{linsys}{3}
     3x  &+  &2y  &+  &z  &=  &0  \\
     6x  &+  &4y  &   &   &=  &0  \\
         &   &y   &+  &z  &=  &0
   \end{linsys}
   \grstep{-2\rho_1 +\rho_2}
   \begin{linsys}{3}
      3x  &+  &2y  &+  &z  &=  &0  \\
          &   &    &   &-2z&=  &0  \\
          &   &y   &+  &z  &=  &0
    \end{linsys}
   \grstep{\rho_2 \leftrightarrow\rho_3}
   \begin{linsys}{3}
      3x  &+  &2y  &+  &z  &=  &0  \\
          &   &y   &+  &z  &=  &0  \\
          &   &    &   &-2z&=  &0
    \end{linsys}
\end{equation*}
\end{example}

\begin{example} \label{ex:SolnChemProb}
有些齐次方程组有许多解。
其中一个就是第一节第一个小节首页的化学问题\index{Chemistry problem}。
\begin{align*}
  \begin{linsys}{4}
              7x  &   &   &-  &7z  &   &   &=  &0  \\
              8x  &+  &y  &-  &5z  &-  &2w &=  &0  \\
                  &   &y  &-  &3z  &   &   &=  &0  \\
                  &   &3y &-  &6z  &-  &w  &=  &0
  \end{linsys}
  &\grstep{-(8/7)\rho_1+\rho_2}
  \begin{linsys}{4}
                7x &   &   &-  &7z  &   &   &=  &0  \\
                   &   &y  &+  &3z  &-  &2w &=  &0  \\
                   &   &y  &-  &3z  &   &   &=  &0  \\
                   &   &3y &-  &6z  &-  &w  &=  &0
   \end{linsys}                                        \\
  &\grstep[-3\rho_2+\rho_4]{-\rho_2+\rho_3}
  \begin{linsys}{4}
                7x &   &   &-  &7z  &   &   &=  &0  \\
                   &   &y  &+  &3z  &-  &2w &=  &0  \\
                   &   &   &   &-6z &+  &2w &=  &0  \\
                   &   &   &   &-15z&+  &5w &=  &0
   \end{linsys}                                        \\
  &\grstep{-(5/2)\rho_3+\rho_4}
  \begin{linsys}{4}
                7x &   &   &-  &7z  &   &   &=  &0  \\
                   &   &y  &+  &3z  &-  &2w &=  &0  \\
                   &   &   &   &-6z &+  &2w &=  &0  \\
                   &   &   &   &    &   &0  &=  &0
   \end{linsys}
\end{align*}
解集
\begin{equation*}
  \set{\colvec[r]{1/3 \\ 1 \\ 1/3 \\ 1}w \suchthat w\in\Re}
\end{equation*}
除了零向量之外还有许多向量
（如果把 \( w \) 取为分子的个数，那么只有当 \( w \) 是 $3$ 的非负倍数时，
解才有意义）。
\end{example}

\begin{lemma} \label{le:HomoSltnSpanVecs}
%<*le:HomoSltnSpanVecs>
对任一齐次线性方程组，存在向量 $\vec{\beta}_1$, \ldots, $\vec{\beta}_k$，
使得该方程组的解集为
\begin{equation*}
  \set{c_1\vec{\beta}_1+\cdots+c_k\vec{\beta}_k \suchthat c_1,\ldots,c_k\in\Re}
\end{equation*}
其中 $k$ 是该方程组的阶梯形形式中
自由变量的个数。
%</le:HomoSltnSpanVecs>
\end{lemma}

在证明之前，我们先讨论关于它的几点。
第一，该证明验证的是：给定一个阶梯形的齐次方程组，
我们可以用回代把所有首变量
都用自由变量表示出来。
要明白为什么这足以建立该引理，我们给出一个例子
（我们也将借这次逐步演算来说明
证明中的记号）。

%<*ex:HomoSystemGivesSpanningSet0>
考虑这个阶梯形的齐次方程组。
\begin{equation*}
  \begin{linsys}{5}
     x  &+  &y   &+  &2z  &+  &u  &+  &v  &=  &0  \\
        &   &y   &+  &z  &+  &u  &-  &v  &=  &0  \\
        &   &    &   &   &   &u  &+  &v  &=  &0
  \end{linsys}
\end{equation*}
它有五个变量：$x_1=x$, \ldots{}, $x_5=v$。
做回代时，我们从最底下的行，即方程~$m=3$ 开始。
利用它的首项 $a_{m,\ell_m}x_m=a_{3,4}x_4=1\cdot u$，
我们解出它的首变量 $x_{\ell_3}=x_4=u$
（记号中的“$\ell$”代表“leading”，即“首”）。
我们得到 $u=-v$。
%</ex:HomoSystemGivesSpanningSet0>

现在回代进入迭代阶段。
%<*ex:HomoSystemGivesSpanningSet1>
我们将用变量~$t$ 来数当前行位于底部之上多少行，
于是取 $t=1$，
考察行 $m-t=3-1=2$。
来到这一行时，我们已经把行~$3$ 的首变量
用自由变量表示了出来。
将 $x_{\ell_3}=x_4=u$ 用它关于
自由变量的表达式代入，得到 $y+z+(-v)-v=0$。
然后解出这一行的首变量，得
$x_{\ell_{m-t}}=x_2=y=-z+2v$。

接下来是位于底部之上 $t=2$ 的那一行，即行 $m-t=1$。
至此我们已经把行 $3$ 和行~$2$ 的首变量
用自由变量表示了出来。
代入 $x_{\ell_3}=u$ 与 $x_{\ell_2}=y$，得到
$x+(-z+2v)+2z+(-v)+v=0$。
然后解出首变量关于自由变量的表达式
$x_{\ell_1}=x=-z-2v$。
%</ex:HomoSystemGivesSpanningSet2>

%<*ex:HomoSystemGivesSpanningSet3>
这个例子的要点在于：现在我们已经完成了。
只要把解写成向量记法
\begin{equation*}
  \colvec{x \\ y \\ z \\ u \\ v}
  =\colvec{-1 \\ -1 \\ 1 \\ 0 \\ 0}z
   +\colvec{-2 \\ 2 \\ 0 \\ -1 \\ 1}v
  \qquad \text{其中 $z,v\in\Re$}
\end{equation*}
就能看出，引理陈述中的 $\vec{\beta}_1$ 与 $\vec{\beta}_2$
就是与自由变量 $z$ 和~$v$ 相关联的向量。
%</ex:HomoSystemGivesSpanningSet3>

关于证明的第二点也可以由那个例子看出。
回代沿方程组自下而上推进，用来自较低行的结论
处理下一行。
这提示了证明策略：数学归纳法。\index{mathematical induction}\index{proof techniques!induction}\index{induction}%

归纳法是一种重要而又不显然的证明技术，本书中
将会多次出现。
我们的归纳证明分两步进行：基础步骤
和归纳步骤。
在基础步骤中，我们验证该命题对某个最初的情形成立，
这里即：对阶梯形线性方程组，我们能写出
最底下方程的首变量关于自由变量的表达式。
在归纳步骤中，我们必须建立一个蕴含关系，
即：如果该命题对所有前面的情形都成立，
那么可以推出该命题对当前情形
也成立。
具体到这里，归纳假设是：如果对
下面的各行我们能把首变量
用自由变量表示出来，那么对再上面一行
我们也能把首变量用这些方式表示出来。

这两个步骤合在一起便对所有的行证明了该命题：
由基础步骤，它对最底下的方程成立，
而由归纳步骤，它对最底下方程成立这一事实
表明它对再上面一个方程也成立。
然后，再用一次
归纳步骤便推出它对再上面第三个方程成立，
如此等等。\appendrefs{mathematical induction}%

\begin{proof}
%<*pf:HomoSltnSpanVecs0>
用高斯消元法化成阶梯形。
其中可能出现一些 \( 0=0 \) 方程；我们忽略它们
（如果
方程组只由 \( 0=0 \)~方程构成，那么该引理显然成立，
因为没有首变量）。
% If there are any variables that appear only with with a coefficient of
% $0$ then they are free, and will never lead a row, so we will
% ignore that possibility.
但因为方程组是齐次的，所以没有矛盾方程。

我们将用归纳法验证
每个首变量都能用
自由变量表示出来。
这将结束证明，因为
我们可以把自由变量用作参数，
而诸 $\vec{\beta}$ 就是由那些
自由变量的系数构成的向量。
%</pf:HomoSltnSpanVecs0>

%<*pf:HomoSltnSpanVecs1>
对于基础步骤，考虑最底下的方程，
\begin{equation*}
  a_{m,\ell_m}x_{\ell_m}+a_{m,1+\ell_{m}}x_{1+\ell_{m}}+\cdots+a_{m,n}x_n=0
   % \tag{$*$}
\end{equation*}
其中 \( a_{m,\ell_m}\neq 0 \)。
%</pf:HomoSltnSpanVecs1>
%<*pf:HomoSltnSpanVecs2>
这是最底下一行，所以
首变量之后 % $x_{\ell_{m}+1}$, \ldots{}
的任何变量
必定都是自由变量。
把这些变量移到右端并除以 $a_{m,\ell_m}$
\begin{equation*}
  x_{\ell_m}
  =(-a_{m,1+\ell_{m}}/a_{m,\ell_m})x_{1+\ell_{m}}+\cdots+(-a_{m,n}/a_{m,\ell_m})x_n
\end{equation*}
便将首变量用自由变量表示了出来。
%</pf:HomoSltnSpanVecs2>
（一个技术性细节：~如果最底下的方程中
在~$x_{\ell_m}$ 的右侧没有变量，
那么 \( x_{\ell_m}=0 \)。
这也满足我们正在验证的命题，因为，
正如本小节开头所提及的，
它把 \( x_{\ell_m} \) 写成了
若干个自由变量的和，具体说是零个自由变量的和，
这里采用的惯例是：
平凡的和总计为~$0$。）

%<*pf:HomoSltnSpanVecs3>
对于归纳步骤，假设对第 \( m \) 个方程，
以及\text{第 \( (m-1) \) 个}
方程，等等，直到并包括\mbox{第 \( m-(t-1) \) 个}~方程，
我们都能
把该方程的首变量用自由变量表示出来。
也就是说，
假设“我们能将一行的首变量
用自由变量表示出来”这一命题
对最底下的 $t$~行成立，其中 $1\leq t<m$。
我们必须验证：由此该命题
对上面紧邻的那个方程，
即第 \( (m-t) \) 个方程也成立。

对 \( x_{\ell_m} \), \ldots, \( x_{\ell_{m-(t-1)}} \) 中的每一个，
代入它关于自由变量的表达式。
所得结果的首项为
$a_{m-t,\ell_{m-t}}x_{\ell_{m-t}}$，其中 \( a_{m-t,\ell_{m-t}}\neq 0 \)，
而左端的其余部分
是自由变量的组合。
把自由变量移到
右端并除以
\( a_{m-t,\ell_{m-t}} \)，最后就把
\( x_{\ell_{m-t}} \) 表示成了自由变量的表达式。

基础步骤和归纳步骤都已完成，所以由
数学归纳法原理，该命题为真。
%</pf:HomoSltnSpanVecs3>
\end{proof}

这表明，正如在引理与它的证明之间所讨论的，
我们可以用自由变量
来参数化解集。
我们说向量集
$\set{c_1\vec{\beta}_1+\cdots+c_k\vec{\beta}_k \suchthat c_1,\ldots,c_k\in\Re}$
\definend{由}集合
\( \set{\smash{\vec{\beta}_1},\ldots,\smash{\vec{\beta}_k}} \)
\definend{生成}\index{generated by}\index{generated}，
或者说是~\definend{由它张成的}\index{spanned by}\index{span}。

% The proof mentions a tricky point.
% We follow the convention that the sum of an empty set of vectors is the
% zero vector.
% In particular, we needs this where
% a homogeneous system has a unique solution because it
% takes the \( c \)'s to be the free variables
% and if there is a unique solution then there are no free variables.

为了完成对 \nearbytheorem{th:GenEqPartPlusHomo} 的证明，
下一个引理考虑解集描述中的
特解部分。

\begin{lemma}  \label{th:GenEqPartHomo}
%<*th:GenEqPartHomo>
对一个线性方程组以及它的任意一个特解 $\vec{p}\/$，
其解集等于
% \begin{equation*}
$
  \set{\vec{p}+\vec{h} \suchthat \text{ \( \vec{h} \) 满足
                                相应的齐次方程组}     }$。
%\end{equation*}
%</th:GenEqPartHomo>
\end{lemma}

% So fixing any particular solution gives the above
% description of the solution set.

\begin{proof}
我们将证明两个方向的集合包含：
方程组的任何解都在上述集合中，而且
上述集合中的任何向量都是该
方程组的解。\appendrefs{set equality}

%<*pf:GenEqPartHomo0>
先证第一个方向的包含：如果一个向量满足该方程组，
那么它在上述集合中。
设 \( \vec{s} \) 满足该方程组。
那么 \( \vec{s}-\vec{p} \) 满足相应的
齐次方程组，因为对每个方程下标 \( i \)，
\begin{multline*}
  a_{i,1}(s_1-p_1)+\cdots+a_{i,n}(s_n-p_n) \\
  =(a_{i,1}s_1+\cdots+a_{i,n}s_n)
  -(a_{i,1}p_1+\cdots+a_{i,n}p_n)
  =d_i-d_i
  =0
\end{multline*}
其中 \( p_j \) 与 \( s_j \) 分别是
\( \vec{p} \) 与 \( \vec{s} \) 的第 \( j \) 个分量。
把 \( \vec{s}-\vec{p} \) 记作 \( \vec{h} \)，
便将 \( \vec{s} \) 表成了所需的 \( \vec{p}+\vec{h} \) 形式。
%</pf:GenEqPartHomo0>

%<*pf:GenEqPartHomo1>
再证另一个方向的包含：取形如 $\vec{p}+\vec{h}$ 的向量，
其中 \( \vec{p} \) 满足该方程组，\( \vec{h} \) 满足
相应的齐次方程组，注意到 $\vec{p}+\vec{h}$
满足所给的方程组，因为对任一方程下标~$i$，
\begin{multline*}
  a_{i,1}(p_1+h_1)+\cdots+a_{i,n}(p_n+h_n)  \\
  =(a_{i,1}p_1+\cdots+a_{i,n}p_n)
   +(a_{i,1}h_1+\cdots+a_{i,n}h_n)
  =d_i+0
  =d_i
\end{multline*}
其中如前所述，\( p_j \) 与 \( h_j \) 分别是
\( \vec{p} \) 与 \( \vec{h} \) 的第 \( j \) 个分量。
%</pf:GenEqPartHomo1>
\end{proof}

这两个引理合起来便证明了 \nearbytheorem{th:GenEqPartPlusHomo}。
请用这句口诀记住该定理：
“\( \text{通解} = \text{特解} + \text{齐次解} \)”。

\begin{example} \label{ex:IllusGenEqPartHomo}
下面这个方程组例示了 \nearbytheorem{th:GenEqPartPlusHomo}。
\begin{equation*}
  \begin{linsys}{3}
    x  &+  &2y  &-  &z  &=  &1  \\
    2x &+  &4y  &   &   &=  &2  \\
       &   &y   &-  &3z &=  &0
  \end{linsys}
\end{equation*}
高斯消元法
\begin{equation*}
  \grstep{-2\rho_1+\rho_2}
  \begin{linsys}{3}
    x  &+  &2y  &-  &z  &=  &1  \\
       &   &    &   &2z &=  &0  \\
       &   &y   &-  &3z &=  &0
  \end{linsys}
  \grstep{\rho_2\leftrightarrow\rho_3}
  \begin{linsys}{3}
      x  &+  &2y  &-  &z  &=  &1  \\
         &   &y   &-  &3z &=  &0  \\
         &   &    &   &2z &=  &0
   \end{linsys}
\end{equation*}
表明通解是一个单元素集。
\begin{equation*}
  \set{\colvec[r]{1 \\ 0 \\ 0} }
\end{equation*}
这个单个的向量显然是一个特解。
相应的齐次方程组经同样的行变换化简
\begin{equation*}
  \begin{linsys}{3}
    x  &+  &2y  &-  &z  &=  &0  \\
    2x &+  &4y  &   &   &=  &0  \\
       &   &y   &-  &3z &=  &0
  \end{linsys}
  \grstep{-2\rho_1+\rho_2}
  \repeatedgrstep{\rho_2\swap\rho_3}
  \begin{linsys}{3}
      x  &+  &2y  &-  &z  &=  &0  \\
         &   &y   &-  &3z &=  &0  \\
         &   &    &   &2z &=  &0
   \end{linsys}
\end{equation*}
也给出一个单元素集。
\begin{equation*}
  \set{\colvec[r]{0 \\ 0 \\ 0} }
\end{equation*}
于是，正如本小节开头所讨论的，
在这种只有一个解的情形，通解就是取特解
再加上相应的齐次方程组的唯一解。
\end{example}

\begin{example}
本小节开头还讨论过，通解集为空的情形
也符合 $\text{通解}=\text{特解}+\text{齐次解}$ 的模式。
这个方程组例示了这一点。
\begin{equation*}
  \begin{linsys}{4}
    x  &   &  &+  &z  &+ &w  &=  &-1  \\
   2x  &-  &y &   &   &+ &w  &=  &3   \\
    x  &+  &y &+  &3z &+ &2w &=  &1
  \end{linsys}
  \grstep[-\rho_1+\rho_3]{-2\rho_1+\rho_2}
  \begin{linsys}{4}
    x  &   &  &+  &z  &+ &w  &=  &-1  \\
       &   &-y&-  &2z &- &w  &=  &5   \\
       &   &y &+  &2z &+ &w  &=  &2
   \end{linsys}
\end{equation*}
它无解，因为最后两个方程
相互冲突。
但相应的齐次方程组确实有解，正如所有
齐次方程组一样。
\begin{equation*}
  \begin{linsys}{4}
    x  &   &  &+  &z  &+ &w  &=  &0   \\
   2x  &-  &y &   &   &+ &w  &=  &0   \\
    x  &+  &y &+  &3z &+ &2w &=  &0
  \end{linsys}
  \grstep[-\rho_1+\rho_3]{-2\rho_1+\rho_2}
  \grstep{\rho_2+\rho_3}
  \begin{linsys}{4}
    x  &   &  &+  &z  &+ &w  &=  &0   \\
       &   &-y&-  &2z &- &w  &=  &0   \\
       &   &  &   &   &  &0  &=  &0
  \end{linsys}
\end{equation*}
事实上，这个解集是无穷的。
\begin{equation*}
  \set{\colvec[r]{-1 \\ -2 \\ 1 \\ 0}z+\colvec[r]{-1 \\ -1 \\ 0 \\ 1}w
         \suchthat z,w\in\Re}
\end{equation*}
尽管如此，由于原方程组没有特解，它的
通解集为空\Dash 不存在形如
$\vec{p}+\vec{h}$ 的向量，因为没有 $\vec{p}\:$。
\end{example}

\begin{corollary} \label{co:ThreeKindsSolutionSets}
%<*co:ThreeKindsSolutionSets>
线性方程组的解集要么为空，要么只有一个元素，
要么有无穷多个元素。
%</co:ThreeKindsSolutionSets>
\end{corollary}

\begin{proof}
%<*pf:ThreeKindsSolutionSets0>
这三种情形发生的例子我们都已见过，
所以只需证明没有其他的可能。

首先注意，一个至少有一个非 \( \zero \) 解 $\vec{v}$ 的齐次方程组
有无穷多个解。
这是因为~$\vec{v}$ 的任何标量倍数也满足该齐次
方程组，而且 $\vec{v}$ 的标量倍数组成的集合中有无穷多个向量：
如果 $s,t\in\Re$ 不相等，那么 $s\vec{v}\neq t\vec{v}$，
因为 $s\vec{v}-t\vec{v}=(s-t)\vec{v}$ 非 $\zero$：
$\vec{v}$ 的任何非 $0$ 分量在经非 $0$ 因子 $s-t$ 倍乘后，
都会给出非 $0$ 的值。
%</pf:ThreeKindsSolutionSets0>

%<*pf:ThreeKindsSolutionSets1>
现在应用 \nearbylemma{th:GenEqPartHomo} 得出：解集
\begin{equation*}
  \set{\vec{p}+\vec{h}\suchthat
    \text{\( \vec{h} \) 满足相应的齐次方程组}}
\end{equation*}
要么为空（若不存在特解 \( \vec{p} \)），
要么只有一个元素（若存在 \( \vec{p} \) 且齐次方程组
有唯一解 \( \zero \)），要么是无穷的（若存在
\( \vec{p} \) 且齐次方程组有非 $\zero$ 解，
从而由上一段可知它有无穷多个解）。
%</pf:ThreeKindsSolutionSets1>
\end{proof}

下表概括了影响通解大小的各个因素。

\smallskip
%<*table:KindsSolutionSets>
\begin{center} % \small
\begin{tabular}{r@{}c}
  &\hspace*{2.5em}\begin{tabular}{c}
      \textit{齐次方程组的} \\[-.5ex]
      \textit{解的个数}
    \end{tabular}                            \\[1.7ex]
  \begin{tabular}{r@{\hspace*{.5em}}}
     \ \\[.6ex]
     \textit{特解} \\[-.55ex]
     \textit{是否存在}
  \end{tabular}
  &\begin{tabular}{r|c@{\hspace*{1em}}c} % \cline{2-3}
     \multicolumn{1}{c}{\ }
         &\textit{一个}    &\textit{无穷多个}                    \\
     \cline{2-3}
     \textit{有}
        &\rule{0ex}{16pt}\begin{tabular}{@{}c@{}} 唯一 \\[-.5ex] 解 \end{tabular}
        &\begin{tabular}{@{}c@{}} 无穷多个 \\[-.5ex] 解 \end{tabular}
         \\[2ex] % \hline
     \textit{无}
        &\begin{tabular}{@{}c@{}} 无 \\[-.5ex] 解 \end{tabular}
        &\begin{tabular}{@{}c@{}} 无 \\[-.5ex] 解 \end{tabular}
         \\ % \hline
   \end{tabular}
\end{tabular}
\end{center}
%</table:KindsSolutionSets>
\smallskip

表中上方的那一维比较简单。
对一个线性方程组施行高斯消元法时，
忽略右端的常数，只关注
左端的系数，
那么最终要么每个变量都领头某一行，要么
我们发现某个变量不领头任何行，也就是说，
我们发现某个变量是自由变量。
（我们通过考虑相应的齐次方程组来将
“忽略右端的常数”形式化。）

一个值得注意的特殊情形是方程个数与未知数个数相同的方程组。
这样的方程组将有解，且解唯一，
当且仅当它
化简为一个每个变量都领头所在行的阶梯形方程组
（因为变量的个数与行的个数相同），
而这又当且仅当
相应的齐次方程组有唯一解。

\begin{definition} \label{df:Nonsingular}
%<*df:Nonsingular>
如果一个方阵是某个有唯一解的齐次方程组的系数矩阵，
则称它是\definend{非奇异的}\index{nonsingular!matrix}
\index{matrix!nonsingular}。
否则，即当它是某个有无穷多解的齐次方程组的系数矩阵时，
称它是\definend{奇异的}\index{singular!matrix}\index{matrix!singular}。
%</df:Nonsingular>
\end{definition}

\begin{example}
下面两个矩阵中，第一个是非奇异的而第二个是奇异的
\begin{equation*}
  \begin{mat}[r]
    1  &2  \\
    3  &4
  \end{mat}
  \qquad
  \begin{mat}[r]
    1  &2  \\
    3  &6
  \end{mat}
\end{equation*}
因为下面两个齐次方程组中第一个有唯一解，
而第二个有无穷多解。
\begin{equation*}
  \begin{linsys}[b]{2}
    x &+  &2y  &=  &0  \\
   3x &+  &4y  &=  &0
  \end{linsys}
  \qquad
  \begin{linsys}[b]{2}
    x &+  &2y  &=  &0  \\
   3x &+  &6y  &=  &0
  \end{linsys}
\end{equation*}
我们在定义中作出这一区分，是因为方程个数与变量个数相同的方程组
只有两种表现方式，取决于它的系数矩阵
是非奇异的还是奇异的。
当系数矩阵非奇异时，无论右端常数是什么，方程组
都有唯一解：~例如，高斯消元法表明方程组
\begin{equation*}
  \begin{linsys}{2}
    x  &+  &2y  &=  &a \\
    3x &+  &4y  &=  &b
  \end{linsys}
\end{equation*}
有唯一解 $x=b-2a$ 和  $y=(3a-b)/2$。
另一方面，当系数矩阵奇异时，方程组绝不会有唯一解\Dash 它
要么无解，要么有无穷多解，如下面这两个方程组。
\begin{equation*}
  \begin{linsys}[b]{2}
    x  &+  &2y  &=   &1   \\
   3x  &+  &6y  &=   &2
  \end{linsys}
  \qquad
  \begin{linsys}[b]{2}
    x  &+  &2y  &=   &1   \\
   3x  &+  &6y  &=   &3
  \end{linsys}
\end{equation*}
\end{example}

定义中使用“奇异”一词是因为它的意思是
“偏离一般的预期”。
人们常常想当然地以为变量个数与方程个数相同的方程组
会有唯一解。
因此，我们可以认为这个词隐含着
“麻烦的”，或者至少“不理想的”意味。
（“奇异”一词恰好用来指那些绝不会恰有一个解的方程组，
这有讽刺意味，但它是标准术语。）

\begin{example}
来自 \nearbyexample{ex:FirstExHomoSys}、
\nearbyexample{ex:HomoZeroOnlySol}
和 \nearbyexample{ex:IllusGenEqPartHomo} 的方程组，
其相应的齐次方程组都有唯一解。
因此这些矩阵都是非奇异的。
\begin{equation*}
  \begin{mat}[r]
    3  &4  \\
    2  &-1
  \end{mat}
  \qquad
  \begin{mat}[r]
    3  &2   &1  \\
    6  &-4  &0  \\
    0  &1   &1
  \end{mat}
  \qquad
  \begin{mat}[r]
    1  &2  &-1 \\
    2  &4  &0  \\
    0  &1  &-3
  \end{mat}
\end{equation*}
而 \nearbyexample{ex:SolnChemProb} 中的化学问题
是一个有不止一个解的齐次方程组，所以它的矩阵
是奇异的。
\begin{equation*}
  \begin{mat}[r]
    7  &0  &-7 &0  \\
    8  &1  &-5 &-2 \\
    0  &1  &-3 &0  \\
    0  &3  &-6 &-1
  \end{mat}
\end{equation*}
\end{example}

上面的表格有两个维度。
我们已经考虑了上面那一个：~只考虑方程组的左端，
我们就能判断一个给定的线性方程组落入哪一列；
右端的常数在这件事中不起作用。

表格的另一个维度，
即判断是否存在特解，就更棘手了。
考虑这两个左端相同但右端不同的方程组。
\begin{equation*}
  \begin{linsys}[b]{2}
    3x &+ &2y &= &5  \\
    3x &+ &2y &= &5
  \end{linsys}
  \qquad
  \begin{linsys}[b]{2}
    3x &+ &2y &= &5  \\
    3x &+ &2y &= &4
  \end{linsys}
\end{equation*}
第一个有解而第二个没有，所以在这里右端的常数
决定了方程组是否有解。
我们可能猜想线性方程组的左端
决定解的个数，而右端决定是否有解，
但这个猜测并不正确。
比较这两个右端相同但左端不同的方程组。
\begin{equation*}
  \begin{linsys}[b]{2}
    3x &+ &2y &= &5  \\
    4x &+ &2y &= &4
  \end{linsys}
  \qquad
  \begin{linsys}[b]{2}
    3x &+ &2y &= &5  \\
    3x &+ &2y &= &4
  \end{linsys}
\end{equation*}
第一个有解而第二个没有。
因此，仅凭方程组右端的常数
并不能决定解是否存在。
确切地说，这取决于左端与
右端之间的某种相互作用。

为了对这种相互作用有一些直观认识，
考虑这个方程组，其中一个系数未加指定，
记为变量~$c$。
\begin{equation*}
  \begin{linsys}{3}
    x  &+  &2y  &+  &3z  &=  &1  \\
    x  &+  &y   &+  &z   &=  &1  \\
   cx  &+  &3y  &+  &4z  &=  &0
  \end{linsys}
\end{equation*}
如果 \( c=2 \)，那么该方程组无解，因为左端
第三行是前两行的和，而右端不是。
如果 \( c\neq 2 \)，那么该方程组有唯一解（取 \( c=1 \) 试一试）。
要使方程组有解：如果左端系数矩阵的某一行
是其他几行的线性组合，
那么右端该行的常数必须是
同样那几行的常数的同样组合。

关于这种相互作用的更多直观认识来自对线性组合的研究。
这将是第二章的重点；在那之前，
我们要在本章其余部分完成对高斯消元法本身的研究。
